\begin{abstract}

Building decentralized storage system that is open to everyone to provide storage resource can benefit users by decoupling the durability reasoning of the system with the trust to individual operators.
However, the permissionless setup also bring its unique challenge of dealing with unpredictable correlated failures.
In this work, we proposed a proper definition of fault tolerance of decentralized storage systems, as the key measure of the durability of such systems.
This novel fault model precisely captures the inability of modeling or predicting correlated failures \ie failure domains in a decentralized system.

Motivated by the vulnerability of data loss of prior works under this fault model, we designed \sys, a decentralized storage system that can tolerate realistic correlated failures with a correlation-oblivious design.
Unlike previous works, \sys distributes encoded fragments to as much storage nodes as possible, ensuring the data object is distributed into sufficient failure domains with overwhelming probability regardless of the ground truth failure domains.
Moreover, \sys deploys rateless erasure code to enable fragment repairing of large code in a permissionless system without explicit coordination.
We also discuss on confidential sampling as a measurement against adversarial attacks.
We believe \sys can serve as a good starting point of designing sophisticated durable decentralized storage system against unpredictable correlated failures.
\end{abstract}